\section{AMC 8 2025}
        \begin{enumerate}
        	\item (\textit{AMC 8 2025}) The eight-pointed star, shown in the figure below, is a popular quilting pattern. What percent of the entire $4\times4$ grid is covered by the star?


 


                    \begin{figure}[H]
                        \centering
                        \includegraphics[width=0.5\textwidth]{Problem/AMC_8_AMC_2025_1_Problem_Figure_1.png}
                    \end{figure}
                    


 $\textbf{(A)}\ 40 \qquad \textbf{(B)}\ 50 \qquad \textbf{(C)}\ 60 \qquad \textbf{(D)}\ 75 \qquad \textbf{(E)}\ 80$


 
 \textbf{Answer: }B



\textbf{Solution 1}

we count the amount of shaded half triangles, $8$. we half this to get the amount of full squares it equals. then we add $4$ for the four full squares in the middle. this gets us a total of eight full squares, out of a total of $16$. so $8$ / $16$ is simplified to $1$ / $2$ or $\boxed{\textbf{(B)}~50}$ percent of the entire grid.~pereena

 


 Each of the unshaded triangles has base length $2$ and height $1$, so they all have area $\frac{2 \cdot 1}{2} = 1$. Each of the unshaded unit squares has area $1$. The area of the shaded region is equal to the area of the entire grid minus the area of the unshaded region, or $4^2 - 4 \cdot 1 - 4 \cdot 1 = 8$. The star is then $\frac{8}{16} = \frac{1}{2} = \frac{50}{100}$, or $\boxed{\textbf{(B)}~50}$ percent of the entire grid.—-DummyAMC

 



\textbf{Solution 2}

There are $16$ total squares in the diagram and each square has $2$ triangles whose areas are half the area of a unit square. Thus, the total number of triangles in the diagram is $4^2 \cdot 2 = 32$ triangles. There are $16$ shaded triangles in the diagram, so the area of the star is $\dfrac{16}{32} = \frac{1}{2} = \frac{50}{100}$, or $\boxed{\textbf{(B)}~50}$ percent.~Pi\_in\_da\_box

 



\textbf{Solution 3}

There are $4$ squares that are entirely shaded and $4$ squares that have no shading. This cancels them out. The rest of the squares are half-half. Therefore the shaded region is $\boxed{\textbf{(B)}~50}$ percent of the grid.

 



\textbf{Solution 4}

Note that we can move one triangle from each of the four cells in the middle to each of the four corners. This will leave every cell in the grid with one triangle each, and each triangle has an area of half the area of each cell. Thus, our answer must equal to $\frac{1}{2}$, and so our answer is $\boxed{\textbf{(B)}~50}.$~derekwang2048

 



\textbf{Solution 5}

The shaded area is a $2 \times 2$ square in the middle of the figure combined with $8$ small triangles. Since each small triangle is $\frac{1}{2}$ of a unit square, the star's area is equal to the area of $4 + 8 \cdot \frac{1}{2} = 8$ unit squares, which $\boxed{\textbf{(B)}~50}$ percent of the grid.

 -vockey

 



\textbf{Solution 6}

You can move the extra half a squares and make 4 squares so there are 8/16. The answer is $\boxed{\textbf{(B)}~50}$ percent of 16.

 William

 


	\item (\textit{AMC 8 2025}) The table below shows the Ancient Egyptian hieroglyphs that were used to represent different numbers.

                    \begin{figure}[H]
                        \centering
                        \includegraphics[width=0.5\textwidth]{Problem/AMC_8_AMC_2025_2_Problem_Figure_1.PNG}
                    \end{figure}
                    


 For example, the number $32$ was represented by the hieroglyphs $\cap \cap \cap ||$. What number is represented by the following combination of hieroglyphs?


 
                    \begin{figure}[H]
                        \centering
                        \includegraphics[width=0.5\textwidth]{Problem/AMC_8_AMC_2025_2_Problem_Figure_2.PNG}
                    \end{figure}
                    


 $\textbf{(A)}\ 1,423 \qquad \textbf{(B)}\ 10,423 \qquad \textbf{(C)}\ 14,023 \qquad \textbf{(D)}\ 14,203 \qquad \textbf{(E)}\ 14,230$


 
 \textbf{Answer: }B



\textbf{Solution 1}

The first hieroglyph is worth $10,000$, the next 4 are worth $100 \cdot 4 = 400$, the next $2$ are worth $10 \cdot 2 = 20$, and the last $3$ are worth $1 \cdot 3 = 3$. Therefore, the answer is $10,000 + 400 + 20 + 3 = \boxed{\textbf{(B)}\ 10,423}$

 



\textbf{Solution 2}

\( 1 \cdot 10^4 + 0 \cdot 10^3 + 4 \cdot 10^2 + 2 \cdot 10^1 + 3 \cdot 10^0 = 10423 \)

 



\textbf{Solution 3 (fast version of Solution x)}

From a first inspection, you may notice that there is no thousands hieroglyph. The only option with zero in the thousands place is $\boxed{\textbf{(B)}\ 10,423}$

 


	\item (\textit{AMC 8 2025}) \textit{Buffalo Shuffle-o} is a card game in which all the cards are distributed evenly among all players at the start of the game. When Annika and $3$ of her friends play \textit{Buffalo Shuffle-o}, each player is dealt $15$ cards. Suppose $2$ more friends join the next game. How many cards will be dealt to each player?


 $\textbf{(A)}\ 8 \qquad \textbf{(B)}\ 9 \qquad \textbf{(C)}\ 10 \qquad \textbf{(D)}\ 11 \qquad \textbf{(E)}\ 12$


 
 \textbf{Answer: }C



\textbf{Solution 1}

We start with Annika and $3$ of her friends playing, meaning that there are $4$ players. This must mean that there is a total of $4 \cdot 15 = 60$ cards. If $2$ more players joined, there would be $6$ players, and since the cards need to be split evenly, this would mean that each player gets $\frac{60}{6}=\boxed{\text{(C)\ 10}}$ Buffalo Shuffle-O cards each meaning that the final answer is 10.

 i agree~pereena

 


	\item (\textit{AMC 8 2025}) Lucius is counting backward by $7$s. His first three numbers are $100$, $93$, and $86$. What is his $10$th number?


 $\textbf{(A)}\ 30 \qquad \textbf{(B)}\ 37 \qquad \textbf{(C)}\ 42 \qquad \textbf{(D)}\ 44 \qquad \textbf{(E)}\ 47$


 
 \textbf{Answer: }B



\textbf{Solution 1}

We plug $a=100, d=-7$ and $n=10$ into the formula $a+d(n-1)$ for the $n$th term of an arithmetic sequence whose first term is $a$ and common difference is $d$ to get $100-7(10-1) = \boxed{\text{(B)\ 37}}$.

 



\textbf{Solution 2}

Since we want to find the $9$th number Lucius says after he says $100$, $7$ is subtracted from his number $9$ times, so our answer is $100-(9 \cdot 7) = \boxed{\text{(B)\ 37}}$

 



\textbf{Solution 3}

Using \href{https://artofproblemsolving.com/wiki/index.php?title=Brute\_force}{brute force} and counting backward by $7$s, we have $100, 93, 86, 79, 72, 65, 58, 51, 44, \boxed{\text{(B)\ 37}}$.

 Note that this solution is not practical and very time-consuming.

 



\textbf{Solution 4}

This can be thought of as an arithmetic sequence. Knowing that our first term is $100$, we have to add $7$ to get to our  0th term, $107$. Our answer is then $107 - 10 \cdot 7 = \boxed{\text{(B)\ 37}}$.

 



\textbf{Solution 5}

Setting this as a function $f (x)= 100- 7(x-1)$ where x is the xth number, we can plug in 10 for x. Quick arithmetic gives us the answer of $\boxed{\text{(B)\ 37}}$.

 


	\item (\textit{AMC 8 2025}) Betty drives a truck to deliver packages in a neighborhood whose street map is shown below.
Betty starts at the factory (labled $F$) and drives to location $A$, then $B$, then $C$, before returning to $F$. What is the shortest distance, in blocks, she can drive to complete the route?


 
                    \begin{figure}[H]
                        \centering
                        \includegraphics[width=0.5\textwidth]{Problem/AMC_8_AMC_2025_5_Problem_Figure_1.png}
                    \end{figure}
                    


 $\textbf{(A)}\ 20 \qquad \textbf{(B)}\ 22 \qquad \textbf{(C)}\ 24 \qquad \textbf{(D)}\ 26\qquad \textbf{(E)}\ 28$


 
 \textbf{Answer: }C



\textbf{Solution 1}

Each shortest possible path from $A$ to $B$ follows the edges of the rectangle. The following path outlines a path of $\boxed{\textbf{(C)}\ 24}$ units:

 
                    \begin{figure}[H]
                        \centering
                        \includegraphics[width=0.5\textwidth]{Solution/AMC_8_AMC_2025_5_Solution_Figure_1.png}
                    \end{figure}
                    

 



\textbf{Solution 2}

We can find the shortest distance using a line diagonally from one point to the other, creating a sort of slope, then find the sum of rise and run of the slope, which happens to be the shortest distance, repeat this process until you get back to Point $F$, and you should get $2 + 1 + 3 + 7 + 4 + 2 + 1 + 4$, which is equal to $\boxed{\textbf{(C)}\ 24}$.~Imhappy62789

 



\textbf{Solution 3}

To find the shortest distance Betty drives, we apply Taxicab Geometry (also known as the $L_1$ metric). In a grid system, the distance between any two points $(x_1, y_1)$ and $(x_2, y_2)$ is defined as $d = |x_1 - x_2| + |y_1 - y_2|$.

 Let the bottom-left corner of the grid be the origin $(0,0)$. Based on the provided map, we assign coordinates to the given locations:

 $B = (0, 0)$

 $C = (2, 4)$

 $F = (6, 5)$

 $A = (7, 3)$

 
Betty's route follows the sequence $F \to A \to B \to C \to F$. We calculate the Manhattan distance for each leg of the journey:

 F (6, 5) to A (7, 3):  $|6 - 7| + |5 - 3| = 1 + 2 = 3$

 A( 7, 3) to B (0, 0):  $|7 - 0| + |3 - 0| = 7 + 3 = 10$

 B (0, 0) to C (2, 4):  $|0 - 2| + |0 - 4| = 2 + 4 = 6$

 C (2, 4) to F (6, 5):  $|2 - 6| + |4 - 5| = 4 + 1 = 5$

 
The total shortest distance is the sum of these segments:
 \[3 + 10 + 6 + 5 = 24\]

 The shortest distance Betty can drive to complete the route is $\boxed{24}$.

 


	\item (\textit{AMC 8 2025}) Sekou writes down the numbers $15, 16, 17, 18, 19.$ After he erases one of his numbers, the sum of the remaining four numbers is a multiple of $4.$ Which number did he erase? 


 $\textbf{(A)}\ 15\qquad \textbf{(B)}\ 16\qquad \textbf{(C)}\ 17\qquad \textbf{(D)}\ 18\qquad \textbf{(E)}\ 19$


 
 \textbf{Answer: }C



\textbf{Solution 1}

The sum of all five numbers is $85. Since$85 is $1 more than a multiple of$4, the number being subtracted must be $1 more than a multiple of$$. Thus, the answer is$\boxed{\text{(C)}~17}$. ~Gavin\_Deng

 



\textbf{Solution 2 - Creativity 💡}

Notice the $10$'s place as you erase one number; you get $40$, which is a multiple of $4$, so you can ignore the $10$'s place in the numbers provided altogether because it will always be a multiple of 40. This means that the only relevant numbers are the $1$'s digit ones. If you add the numbers now, you get $5+6+7+8+9=35.$ The closest multiple of $4$ underneath $35$ which can be achieved as a result of subtracting either $5$, $6$, $7$, $8$, or $9$ is $28$, which is achieved by subtracting $7$, therefore the answer is $\boxed{\textbf{(C)}~17}$.

 



\textbf{Solution 3}

The sum of the residues of these numbers modulo $4$ is $-1+0+1+2+3=5 \equiv 1 \pmod 4$. Hence, the number being subtracted must be congruent to $1$ modulo $4$. The only such answer is $\boxed{\textbf{(C)}~17}$.

 \(\textbf{Note}\): To those who are confused about why, for $15$, it said the residue modulo $4$ is $-1$  instead of $3$ (because of $12$), look here. $-1 \equiv 3 \pmod 4$. The same principle above applies even if you replace $-1$ with $3$; two numbers are in the same residue class if their difference is a multiple of the modulus ( \(3-(-1)=4\), which is \(1\times 4\) ). 

 



\textbf{Solution 4}

The sum of all the numbers here is equivalent to $17 \cdot 5$. Subtracting 17 gives us $17 \cdot 4$, which is clearly a multiple of 4. Therefore, the answer is $\boxed{\textbf{(C)}~17}$.~themathkidwoohooooo

 



\textbf{Solution 5}

Since the sum of $15 + \dots + 19$ is odd, we immediately exclude B and D. We further note that if A is true, then E is true. Hence, the answer is C.  

 



\textbf{Solution 6}

We can use the fact that $(a + b) \mod n = [(a \mod n) + (b \mod n)] \mod n$. Notice that $15, 16, 17, 18, 19$ dividing by 4 have remainders $3, 0, 1, 2, 3$. Their sum is 9. It is easy to see that $9-1=8$ is divisible by 4 and so $\boxed{\textbf{(C)}~17}$ is the correct answer.

 



\textbf{Solution 7}

If we use brute force and add $15+16+17+18+19$, we get $85$. When we subtract every number from $85$, we get $70$, $69$, $68$, $67$, $66$, we see that $68$ is divisible by $4$, so $\boxed{\textbf{(C)}~17}$ is the correct answer.

 


	\item (\textit{AMC 8 2025}) On the most recent exam on Prof. Xochi's class,  


 
                \begin{itemize}
                    \item $5$ students earned a score of at least $95\%$
\item $13$ students earned a score of at least $90\%$
\item $27$ students earned a score of at least $85\%$
\item $50$ students earned a score of at least $80\%$
                \end{itemize}


                How many students earned a score of at least $80\%$ and less than $90\%$?


 $\textbf{(A)}\ 8\qquad \textbf{(B)}\ 14\qquad \textbf{(C)}\ 22\qquad \textbf{(D)}\ 37\qquad \textbf{(E)}\ 45$


 
 \textbf{Answer: }D



\textbf{Solution 1}

$50$ people scored at least $80\%$, and out of these $50$ people, $13$ of them earned a score that was not less than $90\%$, so the number of people that scored in between at least $80\%$ and less than $90\%$ is $50-13 = \boxed{\text{(D)\ 37}}$.

 



\textbf{Solution 2}

Let $a$ denote the number of people who had a score of at least $85$, but less than $90$, and let $b$ denote the number of people who had a score of at least $80$ but less than $85$. Our answer is equal to $a+b$. We find $a = 27 - 13 = 14$, while $b = 50 - 27 = 23$. Thus, the answer is $23 + 14 = \boxed{\text{(D)\ 37}}$. 

 -vockey

 


	\item (\textit{AMC 8 2025}) Isaiah cuts open a cardboard cube along some of its edges to form the flat shape shown on the right, which has an area of $18$ square centimeters. What is the volume of the cube in cubic centimeters?


 
                    \begin{figure}[H]
                        \centering
                        \includegraphics[width=0.5\textwidth]{Problem/AMC_8_AMC_2025_8_Problem_Figure_1.PNG}
                    \end{figure}
                    


 $\textbf{(A)}~3\sqrt{3}\qquad\textbf{(B)}~6\qquad\textbf{(C)}~9\qquad\textbf{(D)}~6\sqrt{3}\qquad\textbf{(E)}~9\sqrt{3}$


 
 \textbf{Answer: }A



\textbf{Solution 1}

Each of the $6$ faces of the cube have equal area, so the area of each face is equal to $\frac{18}{6} = 3$, making the side length $\sqrt3$. From this, we can see that the volume of the cube is $\sqrt{3}^3 = \boxed{\textbf{(A)}~3\sqrt{3}}$

 



\textbf{Solution 2}

Let the side length of the cube shown be equal to $s$ centimeters. The surface area of this cube is equal to the area of the net of the cube, which is equal to $18$ square centimeters. The surface area of this cube is also $6s^2$ square centimeters, so we have
 \[6s^2 = 18 \implies s^2 = \frac{18}{6} = 3 \implies s = \sqrt{3}.\]However, we aren't done. We have found that the side length of the cube is $\sqrt{3} cm$, but the question asks for the volume of the cube, which is equal to $s^3 = \sqrt{3}^3 = \boxed{\textbf{(A)}~3\sqrt{3}}$ cubic centimeters.

 


	\item (\textit{AMC 8 2025}) Ningli looks at the $6$ pairs of numbers directly across from each other on a clock. She takes the average of each pair of numbers. What is the average of the resulting $6$ numbers?


 
                    \begin{figure}[H]
                        \centering
                        \includegraphics[width=0.5\textwidth]{Problem/AMC_8_AMC_2025_9_Problem_Figure_1.png}
                    \end{figure}
                    


 $\textbf{(A)}\ 5\qquad \textbf{(B)}\ 6.5\qquad \textbf{(C)}\ 8\qquad \textbf{(D)}\ 9.5 \qquad \textbf{(E)}\ 12$


 
 \textbf{Answer: }B



\textbf{Solution 1}

This solution uses \href{https://letstalkscience.ca/educational-resources/backgrounders/gauss-summation}{Gauss Summation}.

 Our answer is

 \begin{align}\frac{\frac{1+7}{2} + \frac{2+8}{2} + \cdots + \frac{6+12}{2}}{6} &= \frac{\frac{1}{2}((1+7)+(2+8)+\cdots+(6+12))}{6} \\ \\&= \frac{1+2+3+4+5+6+7+8+9+10+11+12}{2 \cdot 6} \\ \\&= \frac{\frac{12 \cdot 13}{2}}{2 \cdot 6} \\ \\&= \frac{78}{12} \\ \\&= \boxed{\textbf{(B)}~6.5}\end{align}

 



\textbf{Solution 2}

We proceed with \href{https://artofproblemsolving.com/wiki/index.php?title=Brute\_force}{brute force}. All of the pairs of opposite numbers on the clock are $(12,6)$, $(1,7)$, $(2,8)$, $(3,9)$, $(4,10)$, $(5,11)$, where order doesn't matter. The averages of each of these pairs are $9, 4, 5, 6, 7,$ and $8$ respectively, and the average these numbers is $\frac{9+4+5+6+7+8}{6}=\boxed{\textbf{(B)}~6.5}$

 



\textbf{Solution 3 (most efficient)}

The problem is asking for the average of all $12$ numbers. To find the average of all $12$ numbers, you find the sum of all the integers from $1$ to $12$ which is $78$, and divide it by $12$ because there are 12 terms. Therefore, the answer is $\frac{78}{12}=\boxed{\textbf{(B)}~6.5}$.

 You could also see that the average of $2$ and $8$ is $5$, only one more that its consecutive pair, $1$ and $7$ which is $4$. Assuming that this is the pattern, we can find the average $4$ to $9$ (6 consecutive integers starting with $4$), and we that find our answer is $\boxed{\textbf{(B)}~6.5}$.

 



\textbf{Solution 4 (easiest to do in your head)}

You can see that the average of each pair increases by $1$ as you go around the circle; $\frac{1+7}{2}=4$, $\frac{2+8}{2}=5$, and etcetera up to $\frac{6+12}{2}=9$. Therefore, the average of all the pairs is simply the average of the lowest and highest pairs since the average of each pair increases by the same amount as you go around the circle. Hence, the answer is the average of 4 and 9: $\frac{4+9}{2}=\boxed{\textbf{(B)}~6.5}$.

 


	\item (\textit{AMC 8 2025}) In the figure below, $ABCD$ is a rectangle with sides of length $AB = 5$ inches and $AD$ = $3$ inches. Rectangle $ABCD$ is rotated $90^\circ$ clockwise around the midpoint of side $DC$ to give a second rectangle. What is the total area, in square inches, covered by the two overlapping rectangles?


 
                    \begin{figure}[H]
                        \centering
                        \includegraphics[width=0.5\textwidth]{Problem/AMC_8_AMC_2025_10_Problem_Figure_1.png}
                    \end{figure}
                    


 $\textbf{(A)}\ 21 \qquad \textbf{(B)}\ 22.25 \qquad \textbf{(C)}\ 23 \qquad \textbf{(D)}\ 23.75 \qquad \textbf{(E)}\ 25$


 
 \textbf{Answer: }D



\textbf{Solution 1}

The area of each rectangle is $5 \cdot 3 = 15$. Then the sum of the areas of the two regions is the sum of the areas of the two rectangles, minus the area of their overlap. To find the area of the overlap, we note that the region of overlap is a square, each of whose sides have length $2.5$ (as they are formed by the midpoint of one of the long sides, the vertex, and also, since it is rotated 90 degrees). Then the answer is $15+15-2.5^2=\boxed{\textbf{(D)}~23.75}$. 

 


	\item (\textit{AMC 8 2025}) A $\textit{tetromino}$ consists of four squares connected along their edges. There are five possible tetromino shapes, $I$, $O$, $L$, $T$, and $S$, shown below, which can be rotated or flipped over. Three tetrominoes are used to completely cover a $3\times4$ rectangle. At least one of the tiles is an $S$ tile. What are the other two tiles?


 
                    \begin{figure}[H]
                        \centering
                        \includegraphics[width=0.5\textwidth]{Problem/AMC_8_AMC_2025_11_Problem_Figure_1.png}
                    \end{figure}
                    


 $\textbf{(A)}I$ and $L\qquad \textbf{(B)} I$ and $T\qquad \textbf{(C)} L$ and $L\qquad \textbf{(D)}L$ and $S\qquad \textbf{(E)}O$ and $T$


 
 \textbf{Answer: }C



\textbf{Solution 1}

The $3\times4$ rectangle allows for $7$ possible places to put the S piece, with each possible placement having an inverted version. One of the cases looks like this:
                    \begin{figure}[H]
                        \centering
                        \includegraphics[width=0.5\textwidth]{Solution/AMC_8_AMC_2025_11_Solution_Figure_1.png}
                    \end{figure}
                    As you can see, there is a hole in the top left corner of the board, which would be impossible to fill using the tetrominos. There are three cases in which a hole isn't created; the S lies flat in the bottom left corner, it lies flat in the top right corner, or it stands upright in the center. All three tilings are shown below.
                    \begin{figure}[H]
                        \centering
                        \includegraphics[width=0.5\textwidth]{Solution/AMC_8_AMC_2025_11_Solution_Figure_2.png}
                    \end{figure}
                    

 
                    \begin{figure}[H]
                        \centering
                        \includegraphics[width=0.5\textwidth]{Solution/AMC_8_AMC_2025_11_Solution_Figure_3.png}
                    \end{figure}
                    

 
                    \begin{figure}[H]
                        \centering
                        \includegraphics[width=0.5\textwidth]{Solution/AMC_8_AMC_2025_11_Solution_Figure_4.png}
                    \end{figure}
                    For each of the inverted cases, the L pieces can be inverted along with the S piece. Because the only cases that fill the rectangle after the S is placed are the ones that use two L pieces, the answer must be $\boxed{\textbf{(C)}~L \ and \ L}$.~bubby617

 


	\item (\textit{AMC 8 2025}) The region shown below consists of 24 squares, each with side length 1 centimeter. What is the area, in square centimeters, of the largest circle that can fit inside the region, possibly touching the boundaries?


 


                    \begin{figure}[H]
                        \centering
                        \includegraphics[width=0.5\textwidth]{Problem/AMC_8_AMC_2025_12_Problem_Figure_1.png}
                    \end{figure}
                    
$\textbf{(A)}\ 3\pi\qquad \textbf{(B)}\ 4\pi\qquad \textbf{(C)}\ 5\pi\qquad \textbf{(D)}\ 6\pi\qquad \textbf{(E)}\ 8\pi$


 
 \textbf{Answer: }C



\textbf{Solution 1}

The largest circle that can fit inside the figure has its center in the middle of the figure and will be tangent to the figure in $8$ points. By the Pythagorean Theorem, the distance from the center to one of these $8$ points is $\sqrt{2^2 + 1^2} = \sqrt5$, so the area of this circle is \(\pi \sqrt{5}^2 = \boxed{\textbf{(C) } 5\pi}.\)

 



\textbf{Solution 2}

Draw the circle in the grid and analyze the radius. Its radius is a little more than 2 but a lot less than 2.5, so the area is a little more than 4π.  So, the area of the circle is $\boxed{\textbf{(C) } 5\pi}$ with a radius of approximately 2.23.- leafy

 


	\item (\textit{AMC 8 2025}) Each of the even numbers $2, 4, 6, \ldots, 50$ is divided by $7$. The remainders are recorded. Which histogram displays the number of times each remainder occurs?

                    \begin{figure}[H]
                        \centering
                        \includegraphics[width=0.5\textwidth]{Problem/AMC_8_AMC_2025_13_Problem_Figure_1.png}
                    \end{figure}
                    


 
 \textbf{Answer: }A



\textbf{Solution 1}

Let's take the numbers 2 through 14 (evens). The remainders will be 2, 4, 6, 1, 3, 5, and 0. This sequence keeps repeating itself over and over. We can take floor(50/14) = 3, so after the number 42, every remainder has been achieved 3 times. However, since 44, 46, 48, and 50 are left, the remainders of those will be 2, 4, 6, and 1 respectively. The only histogram in which those 4 numbers are set at 4 is histogram $\boxed{\textbf{(A)}}$.

 



\textbf{Solution 2}

Writing down all the remainders gives us 

 
 \[2, 4, 6, 1, 3, 5, 0, 2, 4, 6, 1, 3, 5, 0, 2, 4, 6, 1, 3, 5, 0, 2, 4, 6, 1.\]

 In this list, there are $3$ numbers with remainder $0$, $4$ numbers with remainder $1$, $4$ numbers with remainder $2$, $3$ numbers with remainder $3$, $4$ numbers with remainder $4$, $3$ numbers with remainder $5$, and $4$ numbers with remainder $6$. Manually computation of every single term can be avoided by recognizing the pattern alternates from $0, 2, 4, 6$ to $1, 3, 5$ and there are $25$ terms. The only histogram that matches this is $\boxed{\textbf{(A)}}$.

 


 



\textbf{Solution 3}

First, we find all of the multiples of $7$ that are even, and are therefore, in this list. Knowing that $7$ is odd, and that odd+odd=even, we can find all of the even multiples of $7$ by simply finding all of the multiples of $14$ that fit on this list. Doing this, we end up with

 
 \[14, 28, 42.\]

 Now, we can clearly see that there are only $3$ numbers in this list with $0$ as a remainder. This disproves $2$ of our $5$ answers immediately. Since our remaining answers are identical until we reach $3$ as a remainder, we can skip right to there. Now, we need to find all even numbers that leave a remainder of $3$ when divided by $7$. To do this, we add $3$ to all ODD multiples of $7$ because odd+odd=even. This gives us

 
 \[10, 24, 38.\]

 Note that we don't add $3$ to $49$ because that exceeds $50$. This shows us that there are only $3$ numbers on this list that have a remainder of $3$ when divided by $7$, disproving histograms B and C, which say that there are $4$ of these numbers. This reveals the answer to be $\boxed{\textbf{(A)}}$.

 


	\item (\textit{AMC 8 2025}) A number $N$ is inserted into the list $2, 6, 7, 7, 28$. The mean is now twice as great as the median. What is $N$?


 $\textbf{(A)}\ 7\qquad \textbf{(B)}\ 14\qquad \textbf{(C)}\ 20\qquad \textbf{(D)}\ 28\qquad \textbf{(E)}\ 34$


 
 \textbf{Answer: }E



\textbf{Solution 1}

We see the \href{https://artofproblemsolving.com/wiki/index.php?title=Median}{median} is $7$. Since the mean is two times as great as the median, the mean must be $7 \times 2 = 14$. The sum of the numbers without $N$ is $2 + 6 + 7 + 7 + 28 = 8 + 14 + 28 = 50$. Using the sum, we then have:

 $\frac{N + 50}{6} = 14$. 

 Multiplying both sides by $6$, we get:

 $N + 50 = 14 \times 6 = 84$

 $N = 84 - 50 = \boxed{\text{(E)\ 34}}$

 The median of the list is $7$, so the mean of the new list will be $7 \cdot 2 = 14$. Since there are $6$ numbers in the new list, the sum of the $6$ numbers will be $14 \cdot 6 = 84$. Therefore, $2+6+7+7+28+N = 84 \implies N = \boxed{\text{(E)\ 34}}$

 



\textbf{Solution 2}

Since the average right now is $10$, and the median is $7$, we see that $N$ must be larger than $10$, which means that the median of the 6 resulting numbers should be $7$, making the mean of these $14$. We can do $2 + 6 + 7 + 7 + 28 + N = 14 * 6 = 84$. $50 + N = 84$, so $N$ = $\boxed{\text{(E)\ 34}}$

 i feel this explanation doesn’t make much sense ~anonymous

 $\text{This explanation could be better explained ~anonymous}$

 



\textbf{Solution 3}

We try out every option by inserting each number into the list. After trying out each number, we get $\boxed{\text{(E)\ 34}}$

 Note that this is very time-consuming and it is not the most practical solution.

 



\textbf{Solution 4}

We could use answer choices to solve this problem. The sum of the $5$ numbers is $50$. If you add $7$ to the list, $57$ is not divisible by $6$, therefore it will not work. Same thing applies to $14$ and $20$. The only possible choices left are $28$ and $34$. Now you check $28$. You see that $28$ doesn't work because $(28+50) \div 6 = 13$ and $13$ is not twice of the median, which is still $7$. Therefore, only choice left is $\boxed{\text{(E)\ 34}}$

 



\textbf{Solution 5}

We see that in the list, both $7$s are $3$rd and $4$th. So, no matter what number we plug in, we will have two $7$s in $3$rd and $4$th. So, the median will always be $7$. That means we have to find the option that makes the mean $14$. If the mean is $14$, the sum of all the six numbers is $84$. The sum of the five numbers we are given is $50$, so the answer must be $\boxed{\text{(E)\ 34}}$

 


	\item (\textit{AMC 8 2025}) Kei draws a $6$-by-$6$ grid. He colors $13$ of the unit squares silver and the remaining squares gold. Kei then folds the grid in half vertically, forming pairs of overlapping unit squares. Let $m$ and $M$ equal the least and greatest possible number of gold-on-gold pairs, respectively. What is the value of $m+M$?


 
                    \begin{figure}[H]
                        \centering
                        \includegraphics[width=0.5\textwidth]{Problem/AMC_8_AMC_2025_15_Problem_Figure_1.png}
                    \end{figure}
                    


 $\textbf{(A)}\ 12\qquad \textbf{(B)}\ 14\qquad \textbf{(C)}\ 16\qquad \textbf{(D)}\ 18 \qquad \textbf{(E)}\ 20$


 
 \textbf{Answer: }C



\textbf{Solution 1}

We can see that the least number of gold-on-gold pairs will be obtained when the $13$ silver squares are placed on the two sides so that they don't overlap when folded over (because then it will minimize the number of gold-on-golds). We can see that if we split them up $6$ and $7$ on both sides, and then fold it, the number of gold-on-golds will be $18-13 = 5$.

 The maximum number of gold-on-golds will be achieved when the silver squares overlap when folded over, which will increase the number of gold-on-golds. If we align 6 silver squares with each other on each side, and put the last one somewhere else, we get the maximum is $18 - 7 = 11$. Therefore, the answer is $11+5=\boxed{\textbf{(C)}~16}$. 

 


	\item (\textit{AMC 8 2025}) Five distinct integers from $1$ to $10$ are chosen, and five distinct integers from $11$ to $20$ are chosen. No two numbers differ by exactly $10$. What is the sum of the ten chosen numbers?


 

$\textbf{(A)}\ 95\qquad \textbf{(B)}\ 100\qquad \textbf{(C)}\ 105\qquad \textbf{(D)}\ 110\qquad \textbf{(E)}\ 115$


 
 \textbf{Answer: }C



\textbf{Solution}

Note that \textbf{for no two numbers to differ by $10$, every number chosen must have a different units digit}. To make computations easier, we can choose $(1, 2, 3, 4, 5)$ from the first group and $(16, 17, 18, 19, 20)$ from the second group. Then the sum evaluates to $1+2+3+4+5+16+17+18+19+20  (rearrange    1+19+2+18...) = \boxed{\text{(C)\ 105}}$.

 \textbf{Another Way To Compute} For $1+2+3+4+5+16+17+18+19+20$, we can add the first term and the last term, which is $21$. If we add the second term and the second-to-last term it is also $21$. There are $5$ pairs that sum to $21$, so the answer is $21 \times 5$ which equals $\boxed{\text{(C)\ 105}}$. 

 - leafy

 \textbf{Similar solution} One efficient method is to quickly add $(1, 2, 3, 4, 5, 6, 7, 8, 9, 10)$, which is $55$. Then because you took $50$ in total away from $(16, 17, 18, 19, 20)$, you add $50$. $55+50= \boxed{\text{(C)\ 105}}$.

 


 



\textbf{Solution 4 (effecient)}

To solve this problem, I started with the easiest/smallest case possible. In my opinion, that was

 
$1+2+3+4+5+16+17+18+19+20$.

 
I then solved this equation quickly using Little Gauss's method, rearranging that into

 
$1+19+2+18+3+17+4+16+5+20$.

 
I simplified this into

 
$20+20+20+20+25$.

 
Solving this simple equation gives us the answer, which is $\boxed{\text{(C)\ 105}}$.

 



\textbf{Solution 5 (logic: for me it took about 1.5 mins)}

Start with $1$. If you choose it, you cannot choose $11$, so you must now pick $12$. If you choose $12$, then you can't choose $2$, so you must now pick $13$. If you choose $13$, you cannot choose $3$, so you must pick $4$. If you choose $4$, then you cannot choose $14$ so you have to go to $15$. Then you can't choose $5$ so you have to go to $6$. Then you can't choose $16$, so you have to go to $17$. Then you can't choose $7$, so you have to choose $18$. Then you can't choose $8$, so you have to go to $9$. Finally, you cannot choose $19$, so you must go to $20$. This makes the equation $1+12+3+14+5+16+7+18+9+20 = 105$, hence $\boxed{\text{(C)\ 105}}$.

 Sorry if it is a bit confusing its just how I solved it so wanted to share!-N.

 


	\item (\textit{AMC 8 2025}) In the land of Markovia, there are three cities: $A$, $B$, and $C$. There are $100$ people who live in $A$, $120$ who live in $B$, and $160$ who live in $C$. Everyone works in one of the three cities, and a person may work in the same city where they live. In the figure below, an arrow pointing from one city to another is labeled with the fraction of people living in the first city who work in the second city. (For example, $\frac{1}{4}$ of the people who live in $A$ work in $B$.) How many people work in $A$?  


 
                    \begin{figure}[H]
                        \centering
                        \includegraphics[width=0.5\textwidth]{Problem/AMC_8_AMC_2025_17_Problem_Figure_1.png}
                    \end{figure}
                    


 $\textbf{(A)}\ 55\qquad \textbf{(B)}\ 60\qquad \textbf{(C)}\ 85\qquad \textbf{(D)}\ 115\qquad \textbf{(E)}\ 160$


 
 \textbf{Answer: }D



\textbf{Solution 1}

There are $100 \cdot (\frac{1}{4} + \frac{1}{5}) = 100 \cdot \frac{9}{20} = 45$ people who do not work in city $A$ that live in city $A$, meaning that $100 - 45 = 55$ people who live in city $A$ work in city $A$. There are $\frac{1}{3} \cdot 120 = 40$ people who live in city $B$ and work in $A$, as well as $\frac{1}{8} \cdot 160 = 20$ people who live in city $C$ that work in city $A$. Therefore, the answer is $55 + 40 + 20 = \boxed{\textbf{(D)}\ 115}$.

 


 



\textbf{Solution 2}

We could also make an equation. Let's denote the number of people that live in city $A$ as $a$, $B$ as $b$, and $C$ as $c$. If $x$ denotes the number of people working in city $A$, then $x = \frac{1}{8}c + \frac{1}{3}b + a - \frac{1}{4}a - \frac{1}{5}a$. This is because $\frac{1}{8}$ of the people from City $C$ and $\frac{1}{3}$ of the people from City $B$ work in city $A$ as shown in the image. We also know that $\frac{9}{20}$ of the people living in $A$ work in other cities. We are already given the values of variables $a$, $b$, and $c$ as $100$, $120$, and $160$ respectively. Plug the values it into the main equation like this: $x = 160 \cdot \frac{1}{8} + 120 \cdot \frac{1}{3} + 100 - 100 \cdot \frac{1}{4} - 100 \cdot \frac{1}{5}$. We solve for it and get our answer $\boxed{\textbf{(D)}\ 115}$.

 \textbf{Remark} 
            \begin{enumerate}[label=(\roman*)]
                \item <i>Main article: \href{https://artofproblemsolving.com/wiki/index.php?title=Markov\_Chains}{Markov Chains}</i>
            \end{enumerate}
            This model is known as the \href{https://artofproblemsolving.com/wiki/index.php?title=Markov\_Chains}{Markov Chain}, a type of stochastic process that models systems where the next state depends only on the current state, not on the sequence of events that preceded it. This is known as the Markov property (memoryless property). 

 


	\item (\textit{AMC 8 2025}) The region between a circle of radius 1 and an inscribed square is shaded. Another circle has radius $R$. A quarter of the region between this second circle and its inscribed square is shaded. If the two shaded regions have the same area, what is the value of $R$?


 
                    \begin{figure}[H]
                        \centering
                        \includegraphics[width=0.5\textwidth]{Problem/AMC_8_AMC_2025_18_Problem_Figure_1.png}
                    \end{figure}
                    


 $\textbf{(A)}\ \sqrt2\qquad \textbf{(B)}\ 2\qquad \textbf{(C)}\ 2\sqrt2\qquad \textbf{(D)}\ 4\qquad \textbf{(E)}\ 4\sqrt2$


 
 \textbf{Answer: }B



\textbf{Solution 1}

The area of the shaded region in the circle on the left is the area of the circle minus the area of the square, or $\big(\pi-2)$. The shaded area in the circle on the right is $\dfrac{1}{4}$ of the area of the circle minus the area of the square, or $\dfrac{\pi R^2-2R^2}{4}$, which can be factored as $\dfrac{R^2(\pi-2)}{4}$. Since the shaded areas are equal to each other, we have $\pi-2=\dfrac{R^2(\pi-2)}{4}$, which simplifies to $R^2=4$. Taking the square root, we have $R=\boxed{\text{(B)\ 2}}$

 



\textbf{Solution 2 (Trig)}

The side length of the smaller square is $\sqrt{2}$, which follows from the Pythagorean Theorem since the square is inscribed in a circle of radius $1$, making the diagonal equal to $2$. Thus, the area of the square is $(\sqrt{2})^2 = 2$. The circle in which the square is inscribed has radius $1$, so its area is $\pi \cdot 1^2 = \pi$. Therefore, the area of the shaded region (the part of the circle not covered by the square) is $\pi - 2$.

 Now, consider the larger square, which is inscribed in a circle of radius $R$. This square can be divided into four congruent right triangles by drawing lines from the center of the circle to the vertices of the square. Each triangle has legs of length $R$ and includes a $90^\circ$ angle at the center. Using the triangle area formula $\frac{1}{2}ab\sin(C)$ with $a = b = R$ and $C = 90^\circ$, each triangle has area $\frac{1}{2}R^2$. Since there are four such triangles, the area of the square is $4 \cdot \frac{1}{2}R^2 = 2R^2$. The area of the circle is $\pi R^2$, so the area of the shaded region (the part of the circle outside the square) is $\frac{1}{4}(\pi R^2 - 2R^2) = \frac{1}{4}R^2(\pi - 2)$.

 We are told that this shaded region also has area equal to $\pi - 2$, so we set up the equation $\frac{1}{4}R^2(\pi - 2) = \pi - 2$. Assuming $\pi - 2 \neq 0$, we divide both sides by $\pi - 2$, giving $\frac{1}{4}R^2 = 1$, which leads to $R^2 = 4$, and finally $R = 2$. Thus, the answer is $\boxed{\text{(B) } 2}$.

 



\textbf{Solution 3}

We start with the first area. Since the square is inscribed, its diagonal is $2\implies$ its side length is $\sqrt{2}\implies$ its area is $2$, therefore the first area is $\pi-2$. The second area is $\dfrac{R^2\pi-2R^2}{4}$, found in a similar manner. Writing and solving the equation, we have 
 \[\pi-2=\dfrac{R^2\pi-2R^2}{4}\implies4(\pi-2)=R^2(\pi-2)\implies R=2.\]The answer is $\boxed{\text{(B) }2}$ ~Tacos\_are\_yummy\_1

 



\textbf{Solution 4 (Using similarity)}

Since the two figures are similar, and the 4 regions in the first circle are equal in area to one of the four regions in the second circle, the second figure's area is $4$ times the first figure's area. Therefore, the side lengths are multiplied by a factor of $\sqrt{4} = 2$, so the answer is $1 \cdot 2 = \boxed{\textbf{(B)} 2}$.

 



\textbf{Solution 5}

(Really similar to Solution 1, but has a more larger explanation for the answer)First find the area of the smaller circle with the known radius. Since the radius is $1$, the area of the circle is $\pi$, and finding the side length of the square would equal to $\sqrt{2}$. Squaring that gives us the area of the square which is $2$, and the shaded area of the first one is equal to $\pi - 2$. The second area can be shortened to $\frac{R^2\pi}{4} - \frac{\left(\sqrt{2R^2}\right)^2}{4}$, since we only need $\frac{1}{4}$th of the sections. We then make them equal, which makes this equation:

 $\pi - 2 = \frac{R^2\pi}{4} - \frac{\left(\sqrt{2R^2}\right)^2}{4}$

 Simplifying this equation would be $4\pi - 8 = R^2\pi - 2R^2$. We can factor out the $R^2$ which makes $4\pi - 8 = R^2\left(\pi - 2\right)$. Dividing out the $\pi - 2$ would give us $4$ on the left side. $4 = R^2$ which leaves us to choice $\boxed{\text{(B) }2}$.~Imhappy62789

 



\textbf{Solution 6}

We are given two circles, each with an inscribed square. In the first circle, the radius is 1 unit, and the region between the circle and the inscribed square is shaded. In the second circle, one quarter of the region between the circle and the inscribed square is shaded. The shaded areas in both circles are the same. Our goal is to determine the radius $R$ of the second circle.

 \textbf{Step 1: Set up a ratio for the areas of the shaded regions}

 We know that the shaded region in the first circle (with radius 1) and the shaded region in the second circle (with radius $R$) are proportional in terms of the total area between the circle and the inscribed square. The shaded area in the second circle is one quarter of the total area between the circle and the square, while the shaded area in the first circle is the entire region.

 \textbf{Step 2: Relate the areas to the radius}

 Let’s recall that the area between the circle and the inscribed square is proportional to the square of the radius of the circle because the square's side length scales with the radius of the circle.

 So, in the first circle with radius 1:
 \[\text{Area between circle and square} \propto 1^2 = 1\]and the shaded area is proportional to the entire area between the circle and the square, i.e., 1.

 In the second circle with radius $R$:
 \[\text{Area between circle and square} \propto R^2\]and the shaded area is one-quarter of this total area, so it is proportional to $\frac{1}{4} R^2$.

 \textbf{Step 3: Set up the ratio of shaded areas}

 We are told that the shaded areas in both circles are the same, so we set up the following equation:
 \[1 = \frac{1}{4} R^2\]

 \textbf{Step 4: Solve for $R$}

 Multiply both sides of the equation by 4:
 \[4 = R^2\]Taking the square root of both sides:
 \[R = 2\]

 Thus, the radius of the second circle is $\boxed{\text{(B) 2}}$.

 



\textbf{Solution 7}

Notice that if you shaded the rest of the region between the square and the circle, then the area would be $4$ times the area of the shaded region in the smaller circle. Since the radius changes by a factor of $k^2$,  $k^2 = 4$, solving for gets you $2$. Hence, the answer is $\fbox{\textbf{(B)}\ 2}$

 



\textbf{Solution 8 (Very similar to the ones above)}

We want to find a quick and easy formula that can be applied to this (not necessarily a formula, but a rule). We take the area of the circle and subtract it from the area of the square and you will find the area of the $4$ outer shaded regions. This rule is basically $\pi r^2-2r^2=r^2(\pi-2)$. We plug in the radius from the first figure and compute the answer to be $\pi-2$. Since this represents the area of ONE region in the second figure, we have to multiply our result by four which gets us $4\pi-8$ or to make it easier, you can factor the four and make it $4(\pi-2)$. Using our rule, we can set both equations equal to each other. 

 $R^2(\pi-2)=4(\pi-2)$

 $R^2=4$

 $R=\boxed{\textbf{(B)}2}$

 



\textbf{Solution 9 (Do Not Ever Use This)}

Both area formulas, $\pi r^2$ and $s^2$ are quadratic, therefore multiplying radius/side length by $2$ gives $4x$ the area. The shaded region of the larger figure is 4 times the smaller, therefore the answer is $\boxed{\textbf{(B)}2}$. (For anyone wondering the LaTeX for the answer box is putting this thing: "\boxed{\textbf{(AnswerLetter)}AnswerNumber}" between two $ signs.)~themathkidwoohooooo (pls don't use this)

 


	\item (\textit{AMC 8 2025}) Two towns, $A$ and $B$, are connected by a straight road, $15$ miles long. Traveling from town $A$ to town $B$, the speed limit changes every $5$ miles: from $25$ to $40$ to $20$ miles per hour (mph). Two cars, one at town $A$ and one at town $B$, start moving toward each other at the same time. They drive at exactly the speed limit in each portion of the road. How far from town $A$, in miles, will the two cars meet?


 
                    \begin{figure}[H]
                        \centering
                        \includegraphics[width=0.5\textwidth]{Problem/AMC_8_AMC_2025_19_Problem_Figure_1.png}
                    \end{figure}
                    


 $\textbf{(A)}\ 7.75\qquad \textbf{(B)}\ 8\qquad \textbf{(C)}\ 8.25\qquad \textbf{(D)}\ 8.5\qquad \textbf{(E)}\ 8.75$


 
 \textbf{Answer: }D



\textbf{Solution 1}

The first car, moving from town $A$ at $25$ miles per hour, takes $\frac{5}{25} = \frac{1}{5} \text{hours} = 12$ minutes. The second car, traveling a

 


	\item (\textit{AMC 8 2025}) Sarika, Dev, and Rajiv are sharing a large block of cheese. They take turns cutting off half of what
remains and eating it: first Sarika eats half of the cheese, then Dev eats half of the remaining half,
then Rajiv eats half of what remains, then back to Sarika, and so on. They stop when the cheese is
too small to see. About what fraction of the original block of cheese does Sarika eat in total?


 $\textbf{(A)}\ \frac{4}{7}\qquad \textbf{(B)}\ \frac{3}{5}\qquad \textbf{(C)}\ \frac{2}{3}\qquad \textbf{(D)}\ \frac{3}{4}\qquad \textbf{(E)}\ \frac{7}{8}$


 
 \textbf{Answer: }A



\textbf{Solution 1}

Let the total amount of cheese be $1$. We will track the amount of cheese Sarika eats throughout the process.

 \textbf{First Round:} Sarika eats half of the total cheese, so she eats:
 \[\frac{1}{2}.\]\textbf{Second Round:} Dev eats half of what remains after Sarika's turn, which is:
 \[\frac{1}{4}.\]\textbf{Third Round:} Rajiv eats half of the remaining cheese after Dev’s turn, which is:
 \[\frac{1}{8}.\]

 At the end of the first round, the total cheese eaten is:
 \[\frac{1}{2} + \frac{1}{4} + \frac{1}{8} = \frac{7}{8}.\]

 We observe that Sarika’s consumption follows a <b>\href{https://artofproblemsolving.com/wiki/index.php?title=Geometric\_sequence}{geometric sequence}</b>. In the first round, she eats $\frac{1}{2}$, and in subsequent rounds, she eats half of what remains from the previous round. This gives the following series for Sarika’s total consumption:

 
 \[\frac{1}{2} + \frac{1}{16} + \frac{1}{128} + \cdots\]

 This is a geometric series with first term $\frac{1}{2}$ and common ratio $\frac{1}{8}$. The sum $S$ of this infinite geometric series is given by the formula:

 
 \[S = \frac{a}{1 - r},\]where $a$ is the first term and $r$ is the common ratio. Substituting $a = \frac{1}{2}$ and $r = \frac{1}{8}$:

 
 \[S = \frac{\frac{1}{2}}{1 - \frac{1}{8}} = \frac{\frac{1}{2}}{\frac{7}{8}} = \frac{1}{2} \times \frac{8}{7} = \frac{4}{7}.\]

 Thus, Sarika eats $\frac{4}{7}$ of the original block of cheese. The correct answer is:
 \[\boxed{\textbf{(A) } \frac{4}{7}}.\]

 



\textbf{Solution 2 (Estimation)}

Sarika eats $1/2$ of the original cheese, and then because the others eat $1/4$ and $1/8$, she eats $1/16$ next, and then $1/128$, and then so on. Since the values later are going to be too small to make a huge difference, we can use these $3$ values. She ate $(64 + 8 + 1)/128 = 73/128$. We can replace the $73$ with a $72$ for now, so $72/128 = 9/16$, which simplifies to around $56.25$. Since there is a little bit more of the cheese to be accounted for, the amount that she eats will be around $\boxed{\textbf{(A) }\frac{4}{7}}$.

 



\textbf{Solution 3 Speed}

Her first $2$ bites will be $\frac{1}{2}$ and $\frac{1}{16}$. Everything else won't matter that much, so the answer will be closest to $\frac{9}{16}$, $\frac{4}{7}$.

 



\textbf{Solution 4 (Basically Video Solution)}

Let Sarika eat a total of $x$ units of cheese. We can then see that Dev will eat a total of $\frac{x}{2}$ units and Rajiv will eat $\frac{x}{4}$ units. Our desired answer is $\frac{x}{x + \frac{x}{2} + \frac{x}{4}}$, or $\boxed{\textbf{(A) }\frac{4}{7}}$.

 


 
Key Idea: Let $x$ be the fraction eaten by Sarika. Then Dev eats $\frac{x}{2}$ and Rajiv eats $\frac{x}{4}$. Hence $x + \frac{x}{2} + \frac{x}{4} = 1$ so solving for $x$ we get 
 \[x = \boxed{\textbf{(A) }\frac{4}{7}}\].

 


	\item (\textit{AMC 8 2025}) The Konigsberg School has assigned grades 1 through 7 to pods $A$ through $G$, one grade per
pod. Some of the pods are connected by walkways, as shown in the figure below. The school
noticed that each pair of connected pods has been assigned grades differing by 2 or more grade
levels. (For example, grades 1 and 2 will not be in pods directly connected by a walkway.) What is
the sum of the grade levels assigned to pods C, D, E, and F? 


 
                    \begin{figure}[H]
                        \centering
                        \includegraphics[width=0.5\textwidth]{Problem/AMC_8_AMC_2025_21_Problem_Figure_1.png}
                    \end{figure}
                    


 $\textbf{(A)}~12\qquad\textbf{(B)}~13\qquad\textbf{(C)}~14\qquad\textbf{(D)}~15\qquad\textbf{(E)}~16$


 
 \textbf{Answer: }A



\textbf{Solution 1}

The key observation for this solution is to observe that pods $C$ and $F$ both have degree 5; that is, they are connected to five other pods. This implies that $C$ and $F$ must contain grades 1 and 7 in either order (if $C$ or $F$ contained any of grades 2, 3, 4, 5, or 6, then there would only be four possible grades for five pods, a contradiction). It does not matter which grade $C$ and $F$ gets (see Solution 2 below), so we will assume that pod $C$ is assigned grade 1, and pod $F$ is assigned grade 7.

 Next, pod $D$ is the only pod which is not adjacent to pod $F$, so pod $D$ must be assigned grade 6. Similarly, pod $G$ must be assigned grade 2.

 Lastly, we need to assign grades 3, 4, and 5 to pods $A$, $B$, and $E$. Note that $A$ and $B$ are adjacent; therefore, pods $A$ and $B$ must contain grades 3 and 5. We conclude that $E$ must be assigned grade 4. The requested sum is $1+4+7 = \mathbf{(A)\,} 12$.

 



\textbf{Solution 2}

Suppose that we replaced each label $x$ with $8-x$. That is, we replaced the grade $a$ assigned to $A$ with the grade $8-a$, the grade $b$ assigned to $B$ with $8-b$, and so on. We claim that if the initial labeling is valid (i.e., each pair of connected pods has grades differing by 2 or more), then so will the new labeling. This follows as if $|x-y| \ge 2$, then $|(8-x)-(8-y)|=|-x+y| = |x-y| \ge 2$. This shows that if there is an assignment $(a,b,c,d,e,f,g)$ of grades to the pods, where $\{a,b,\ldots,g\} = \{1,2,\ldots,7\}$, then there is a second, valid assignment.

 But because this is an AMC 8 problem and there is only one correct answer, we can conclude that $c+e+f = (8-c)+(8-e)+(8-f)$. Simplifying and solving for $c+e+f$ gives $c+e+f = \mathbf{(A)\,} 12$.

 


	\item (\textit{AMC 8 2025}) A classroom has a row of 35 coat hooks. Paulina likes coats to be equally spaced, so that there is the same number of empty hooks before the first coat, after the last coat, and between every coat and the next one. Suppose there is at least 1 coat and at least 1 empty hook. How many different numbers of coats can satisfy Paulina's pattern?


 
                    \begin{figure}[H]
                        \centering
                        \includegraphics[width=0.5\textwidth]{Problem/AMC_8_AMC_2025_22_Problem_Figure_1.png}
                    \end{figure}
                    $\textbf{(A)}\ 2\qquad \textbf{(B)}\ 4\qquad \textbf{(C)}\ 5\qquad \textbf{(D)}\ 7\qquad \textbf{(E)}\ 9$


 
 \textbf{Answer: }D



\textbf{Solution 1}

Suppose there are $c$ coats on the rack. Notice that there are $c+1$ "gaps" formed by these coats, each of which must have the same number of empty spaces (say, $k$). Then the values $k$ and $c$ must satisfy $c+k(c+1)=35 \implies kc+k+c=35$. We now use \href{https://artofproblemsolving.com/wiki/index.php?title=Simon\%27s\_Favorite\_Factoring\_Trick}{Simon's Favorite Factoring Trick} as follows: 
 \[kc+k+c=35\] 
 \[\implies kc+k+c+1=36\] 
 \[\implies (k+1)(c+1)=36\] Our only restrictions now are that $k>0 \implies k+1 > 1$ and $c>0 \implies c+1>1$. Other than that, each factor pair of $36$ produces a valid solution $(k,c)$, which in turn uniquely determines an arrangement. Since $36$ has $9$ factors, our answer is $9-2=\boxed{\textbf{(D)}~7}$. ~cxsmi

 



\textbf{Solution 2}

Say Paulina placed $n$ coats. That will divide the 35 hooks into $n+1$ spaces and $35-n$ empty hooks. Therefore, 
 \[n+1|35-n\] The values of $n$ that satisfy this are 
 \[n\in{1,2,3,5,8,11,17}\]The answer is $\boxed{\text{(D) }7}$.

 



\textbf{Solution 3}

Let the number of empty hooks before the first coat be denoted as $n$.

 This means that there are $n$ empty hooks between all the coats. In every situation, there are $n + 1$ gaps, since $n$ cannot be zero.

 To solve the problem, subtract $2n$ from 35 (accounting for the $n$ empty hooks at the beginning and $n$ empty hooks at the end). Then subtract 1 (for the first occupied hook). After that, divide the remaining number of hooks by $n + 1$ (the number of hooks between coats). We do this because the first occupied hook is always after $n$ hooks; we just need to check if the remaining hooks can be evenly divided by the number of gaps.

 To clarify: 

 
                \begin{itemize}
                    \item If $n = 1$:
                \end{itemize}


                \begin{align*}35 - 2n &= 33 \\33 - 1 &= 32 \\\frac{32}{n + 1} &= 16 \end{align*}


 Since 16 is an integer, this situation satisfies the conditions.

 
                \begin{itemize}
                    \item If $n = 2$:
                \end{itemize}


                We can skip the the third step, as each time you add 1 to $n$, it decreases by 2 from the previous situation.

 
 \[\frac{30}{n + 1} = 16\]

 This situation also satisfies the conditions.

 Continue this process until the difference exceeds $35 - (2n + 1)$. We find that 7 situations satisfy the conditions. Therefore, the answer is 7.

 \textit{Note: This method is slow but easy to check and understand.}

 



\textbf{Solution 4 (Similar to Solution 1 but without Simon's Favorite Factoring Trick)}

Suppose there are $n$ occupied hooks on the rack. Observe that there are $n+1$ intervals, each of which must have the same number of empty hooks. Let's suppose $m$ is the number of empty hooks plus the occupied hook to its right. Then, we have the equation 
 \[(n+1)m-1=35\] 
 \[\implies (n+1)m=36\] Our restrictions are n > 0 and m > 1. 36 has 9 factors, and two pairs do not meet the restriction. Our answer is then $9-2=\boxed{\textbf{(D)}~7}$.

 



\textbf{Solution 5 (neat trick)}

Note that, for any possible valid arrangement of coats, there is always a set number of empty coat hangers between the coats. This is never 0, since then there would be no empty hangers on either side. So suppose there is $n$ hangers between every coat. Then, there is one coat every n+1 hangers. For any value of $n+1$, divide it into 35, the amount of hangers. If 
 \[35/(n+1)\] has a remainder of $n$, then that would be a valid arrangement for that value of $n+1$. This works for values of 2, 3, 4, 6, 9, 12, and 18. This means there are $\boxed{\textbf{(D)}~7}$ total possible arrangements.

 -Arm123

 


	\item (\textit{AMC 8 2025}) How many four-digit numbers have all three of the following properties?


 (I) The tens and ones digit are both 9.


 (II) The number is 1 less than a perfect square.


 (III) The number is the product of exactly two prime numbers.


 $\textbf{(A)}\ 0 \qquad \textbf{(B)}\ 1 \qquad \textbf{(C)}\ 2 \qquad \textbf{(D)}\ 3 \qquad \textbf{(E)}\ 4$


 
 \textbf{Answer: }B



\textbf{Solution 1}

The "Condition (II)" perfect square must end in "$00$" because $...99+1=...00$ \textit{Condition (I)}. Four-digit perfect squares ending in "$00$" are ${40, 50, 60, 70, 80, 90}$, which when squaring these numbers, gets you $(1600, 2500, 3600, 4900, 6400, 8100)$ respectively.

 \textit{Condition (II)} also says the number is in the form $n^2-1$. By the <b>\href{https://en.wikipedia.org/wiki/Difference\_of\_two\_squares}{Difference of Squares}</b>, $n^2-1 = (n+1)(n-1)$. Hence: 

 
                \begin{itemize}
                    \item $40^2-1 = (39)(41)$\item $50^2-1 = (49)(51)$\item $60^2-1 = (59)(61)$\item $70^2-1 = (69)(71)$\item $80^2-1 = (79)(81)$\item $90^2-1 = (89)(91)$
                \end{itemize}


                On this list, the only number that is the product of $2$ prime numbers (condition $3$) is $60^2-1 = (59)(61)$, so the answer is $\boxed{\text{(B)\ 1}}$. 

 


	\item (\textit{AMC 8 2025}) In trapezoid $ABCD$, angles $B$ and $C$ measure $60^\circ$ and $AB = DC$. The side lengths are all positive integers, and the perimeter of $ABCD$ is $30$ units. How many non-congruent trapezoids satisfy all of these conditions?


 
                    \begin{figure}[H]
                        \centering
                        \includegraphics[width=0.5\textwidth]{Problem/AMC_8_AMC_2025_24_Problem_Figure_1.png}
                    \end{figure}
                    


 $\textbf{(A)}\ 0 \qquad \textbf{(B)}\ 1 \qquad \textbf{(C)}\ 2 \qquad \textbf{(D)}\ 3 \qquad \textbf{(E)}\ 4$


 
 \textbf{Answer: }E



\textbf{Solution 1}

Let $a$ be the length of the shorter base, and let $b$ be the length of the longer one. Note that these two parameters, along with the angle measures and the fact that the trapezoid is isosceles, uniquely determine a trapezoid. We drop perpendiculars down from the endpoints of the top base. Then the length from the foot of this perpendicular to either vertex will be half the difference between the lengths of the two bases, or $\frac{b-a}{2}$. Now, since we have a 30-60-90 triangle and this side length corresponds to the "30" part, the length of the hypotenuse (one of the legs) is $2 \cdot \frac{b-a}{2} = b-a$. Then the perimeter of the trapezoid is $2(b-a)+a+b=3b-a=30$. The only other stipulation for this trapezoid to be valid is that $b>a$ (which was our assumption). We can now easily count the valid pairs $(a,b)$, yielding $(3,11),(6,12),(9,13),(12,14)$. It is clear that proceeding further would cause $a \geq b$, so we have $\boxed{\textbf{(E)}~4}$ valid trapezoids.~cxsmi

 



\textbf{Solution 2}

Let $x$ be the length of $AB$ and $DC$, and let $b$ be the length of the shorter base. Because $\angle B$ and $\angle C = 60^{\circ}$, the length of the longer base is $b + \frac{x}{2} + \frac{x}{2} = b + x$. Therefore, the perimeter is $3x + 2b = 30$. The number of positive integer pairs $(x, b)$ is $(2,12), (4,9), (6,6), (8,3)$, meaning the answer is $\boxed{\textbf{(E)}~4}$.

 way where you don't have to test pairs:

 We can rearrange $3x + 2b = 30$ to $\frac{30-3x}{2} = b$. We want $\frac{30-3x}{2}$ to be a positive integer so b is a positive integer. 30 is a multiple of 3 and $-3x$ is just subtracting multiples of 3's from 30 so every number on the numerator of $\frac{30-3x}{2}$ is a multiple of 3 under 30 ($x$ is also a positive integer so it can't be larger than 30 or not be a multiple of 3). For $30-3x$ to be divisible by 2 it must be even, so we are finding even multiples of 3 under 30 which are 24, 18, 12, and 6. Each number will give us a valid $x$ and a valid $b$ pair giving us 4 solutions.

 



\textbf{Solution 3}

Let $x$ be the length of the legs of the trapezoid. Draw the angle bisectors of the 120 degree angles. Now you have 2 equilateral triangles, and another figure in between them. Let $x+a$ be the length of the shorter base of the trapezoid, and let $2x+a$ be the length of the longer base. Since the perimeter of the trapezoid is 30, $5x+2a = 30$. Since $x>$0, and $a$ can be negative as long as $x+a>0$, you get 4 solutions for $(x,a)$, namely $(2,10), (4,5), (6,0), (8,-5)$. Any solution with $x\ge10$ would lead to $x+a\le0$. Hence, the answer is $\boxed{\textbf{(E)}~4}$

 



\textbf{Solution 4}

Drop altitudes from angle $A$ and angle $D$. Then, two $30-60-90$ triangles and a rectangle are created. Let the hypotenuse of both of the triangles equal $x$ and let side $AD$ equal $y$. Then, the distance from $B$ and $C$ to the feet of the altitudes that are closest to them is $\frac{x}{2}$, and the distance between the feet of the two altitudes would be $y$. The rest is similar to Solution 2~Soupboy0

 



\textbf{Solution 5 (SIMILAR TO SOLUTION 4)}

We give lengths to the base of the $30-60-90$ triangles on both sides as the trapezoid as variable $x$.  We give lengths of Line $AB$ and the other side as variable $y$.  When we add them, we have $6x+2y=30$ from the sides.  If we lext $x=5$, the other sides won't work so it has to be integers less than $5$ which is $4$ which is the answer.- Spacepandamath13

 



\textbf{Solution 6}

We’re given a **trapezoid $ABCD$** with these facts:

 
                \begin{itemize}
                    \item Angles $\angle B = \angle C = 60^\circ$\item $AB = DC$ (the legs are equal)\item All side lengths are **positive integers**\item **Perimeter** = 30 units\item We’re to find how many **non-congruent** trapezoids satisfy all this.
                \end{itemize}


                ---

 \begin{verbatim}
Step 1: Understand the trapezoid
\end{verbatim}

Because $\angle B = \angle C = 60^\circ$, and those angles are **adjacent to the bases**, the trapezoid is **isosceles**.

 Let:

 
                \begin{itemize}
                    \item $AB = DC = x$ (legs—equal)\item $AD = a$, $BC = b$ (bases—not necessarily equal)\item So: **Perimeter = $a + b + 2x = 30$**
                \end{itemize}


                ---

 \begin{verbatim}
Step 2: Use geometry
\end{verbatim}

From the angle info, drop perpendiculars from points $A$ and $D$ to base $BC$. This creates two right triangles with angles of $60^\circ$, so the **height** of the trapezoid is:

 
 \[\text{height} = x \cdot \sin(60^\circ) = x \cdot \frac{\sqrt{3}}{2}\]

 And the **projection** of the leg on the base is:

 
 \[\text{horizontal part} = x \cdot \cos(60^\circ) = x \cdot \frac{1}{2}\]

 So the total difference between the longer base and shorter base is:

 
 \[b - a = 2 \cdot \left(x \cdot \frac{1}{2}\right) = x \quad \Rightarrow \quad b = a + x\]

 ---

 \begin{verbatim}
Step 3: Plug into perimeter
\end{verbatim}


 \[a + b + 2x = 30 \\ a + (a + x) + 2x = 30 \\ 2a + 3x = 30\]

 This is the key equation:

 > 
 \[> \boxed{2a + 3x = 30} >\]

 ---

 \begin{verbatim}
Step 4: Try integer values
\end{verbatim}

We solve for integers $a$ and $x$:

 
 \[2a = 30 - 3x \Rightarrow a = \frac{30 - 3x}{2}\]

 Try small integer values of $x$ so $a$ is a **positive integer**:

 | $x$ | $a = \frac{30 - 3x}{2}$ | Is $a$ integer?  | $b = a + x$ || --- | ----------------------- | ---------------- | ----------- || 1   | $13.5$                  | ✗                |             || 2   | $12$                    | ✔                | $14$        || 3   | $10.5$                  | ✗                |             || 4   | $9$                     | ✔                | $13$        || 5   | $7.5$                   | ✗                |             || 6   | $6$                     | ✔                | $12$        || 7   | $4.5$                   | ✗                |             || 8   | $3$                     | ✔                | $11$        || 9   | $1.5$                   | ✗                |             || 10  | $0$                     | ✗ (not positive) |             |

 So the valid integer combinations are

 
                \begin{itemize}
                    \item $x = 2$, $a = 12$, $b = 14$\item $x = 4$, $a = 9$, $b = 13$\item $x = 6$, $a = 6$, $b = 12$\item $x = 8$, $a = 3$, $b = 11$
                \end{itemize}


                ---

 \begin{verbatim}
Step 5: Eliminate congruent ones
\end{verbatim}

Each trapezoid has side lengths:

 
                \begin{itemize}
                    \item Legs: $x$\item Bases: $a$ and $b = a + x$
                \end{itemize}


                So all side lengths are different in each case:

 1. $12, 14, 2, 2$

 2. $9, 13, 4, 4$

 3. $6, 12, 6, 6$

 4. $3, 11, 8, 8$

 All four have different sets of side lengths → all are **non-congruent**.

 ---

 \begin{verbatim}
(True) Final Answer:
\end{verbatim}


 \[\boxed{4}\]

 -- ChatGPT 4

 -- Edited by JubilantJelly

 



\textbf{Solution 7}

Explanations and Reasoning to Similar Solutions

 
Note that trapezoid $ABCD$ is isosceles with $B$ and $C$ being the bottom 2 vertices. The shape can be dissected into 3 shapes with 2 congruent right triangles separated by a rectangle. This is done by dropping an altitude, a perpendicular point on the base from each of the 2 top vertices, named $E$ and $F$, respectively below $A$ and $D$. The property that angles $\angle$$B$ and $\angle$$C$ measure $60^\circ$ is significant here, forming $30-60-90$ triangles ($\triangle$$ABE$ and $\triangle$$DCF$). Now the assumed larger $^a$ side $BC$ is separated into 3 sections by the 3 shapes. Or $BC = BE + EF + CF$. $^b$  Recall that $\triangle$$ABE$ and $\triangle$$DCF$ are $30-60-90$ triangles, the side length ratios are $1:\sqrt3:2$. Naming the length of the shorter side of the right triangles as $a$, and the length of side $AD$ as $b$, then the trapezoid's side lengths can be expressed in only 2 variables. Now knowing $AD = b$, $DC = 2a$ $^c$, $BC = a + b + a = 2a + b$, and $AC = 2a$, the perimeter is $b + 2a + 2a + b + 2a = 2b + 6a = 30$. Simplifying, $3a + b = 15$, only 4 pairs of values satisfy this condition, namely $(1,12)$, $(2,9)$, $(3,6)$, $(4,3)$.

 Thus there $\boxed{\textbf{(E)}~4}$ possible non-congruent trapezoids. 

 $^a$ The question asks for non-congruent trapezoids, and rotation preserves congruence so  the assumption is always true WLOG.

 $^b$ The shortest side of $\triangle$ABE, the length of AD (shape ADFE is a rectangle so FE = AD), and the shortest side of $\triangle$DCF.

 $^c$ $AB = DC = 2a$, $AB$ and $DC$ are the hypotenuse of triangles $\triangle$ABE and $\triangle$DCF.~moracle

 


 



\textbf{Solution 8}

Let $a = AD$ and $s = AB = DC$. Draw a segment from $A$ to $\overline{BC}$ parallel to $\overline{DC}$, meeting $\overline{BC}$ at point $P$, forming parallelogram $ADCP$. It follows that $AP = DC = s$ and $PC = AD = a$. Isosceles $\triangle ABP$ has base angles measuring $60^\circ$, so it is equilateral and $BP = s$.

 
                    \begin{figure}[H]
                        \centering
                        \includegraphics[width=0.5\textwidth]{Solution/AMC_8_AMC_2025_24_Solution_Figure_1.png}
                    \end{figure}
                    

 Then the perimeter of $ABCD$ is $L = 2a + 3s = 30$. To satisfy this equation with positive integer values, the number $a$ must be divisible by $3$ because $2a = 3(10 - s)$. The only possible values for $a$ are $3, 6, 9,$ and $12$, each value corresponding to a different trapezoid. Greater multiples of $3$ will not lead to a positive integer value for $s$. Therefore there are $\boxed{4}$ different trapezoids that satisfy the conditions.

 The table below lists the side lengths $a$, $s$, and $a + s$ for these $4$ trapezoids, along with their perimeters.

 \begin{array}{c|c|c|c}a & s & a + s & L = 2a + 3s \\\hline3 & 8 & 11 & 2 \cdot 3 + 3 \cdot 8 = 30 \\6 & 6 & 12 & 2 \cdot 6 + 3 \cdot 6 = 30 \\9 & 4 & 13 & 2 \cdot 9 + 3 \cdot 4 = 30 \\12 & 2 & 14 & 2 \cdot 12 + 3 \cdot 2 = 30 \\\end{array}

 



\textbf{Solution 9 (Similar to Solutions 4 and 5)}

We can draw altitudes from A and B and create a rectangle and $2$ $30-60-90$ triangles. Two sides of each triangle are on the trapezoid. We can label each of the hypotenuses of the small $30-60-90$ triangles as $2a$ and each of the short sides of the triangles as $a$. Then, the perimeter of the trapezoid becomes $2(2a + a) +$ a positive integer (we'll call this integer $n$) $= 30$. Simplifying, we get $6a + n = 30$, and since $n$ is not 0, $a$ must be a positive integer less than $5$, which yields $\boxed{4}$ options: $1$, $2$, $3$, and $4$.

 \textbf{solution 10} Extend lines AB and CD until they intersect. Calling that point E, you can find out that triangle EBC is equilateral because the angles at B and C are 60 degrees. Assume that side BC=a, then BE and CE are also equal to a. Assuming that AE=x, then DE must also be x because triangle AED is equilateral. This also means that AD=x. Since CD=CE-DE, we have CD=a-x. Similarly, AB=a-x because AB=CD. This means the perimeter of the trapezoid ABCD is a+(a-x)+(a-x)+x. The problem tells us that the perimeter is 30 so we get the equation 3a-x=30. Adding x and dividing by 3 you get a=10-x/3. Since x is a side length, it must be an integer. That means x can any integer between 0 and 30 exclusive. Side length a must also be an integer, that means x must be a multiple of 3. So x can be 3, 6, 9, 12, 15, 18, 21, 24, 27. From the image given to us, we see that BC=a is greater than AD=x. That means x can't be greater than or equal to 15 because that would cause x to be greater than a. This leaves us with 4 values of x which are 3,6,9,12. That means there are 4 non-congruent trapezoids that satisfy those conditions. This means the answer is $\boxed{\textbf{(E)}~4}$.

 


	\item (\textit{AMC 8 2025}) Makayla finds all the possible ways to draw a path in a $5 \times 5$ diamond-shaped grid. Each path starts at the bottom of the grid and ends at the top, always moving one unit northeast or northwest. She computes the area of the region between each path and the right side of the grid. Two examples are shown in the figures below. What is the sum of the areas determined by all possible paths?


 
                    \begin{figure}[H]
                        \centering
                        \includegraphics[width=0.5\textwidth]{Problem/AMC_8_AMC_2025_25_Problem_Figure_1.png}
                    \end{figure}
                    


 $\textbf{(A)}\ 2520 \qquad \textbf{(B)}\ 3150 \qquad \textbf{(C)}\ 3840 \qquad \textbf{(D)}\ 4730 \qquad \textbf{(E)}\ 5050$


 
 \textbf{Answer: }B



\textbf{Solution 1}

\textbf{Step 1:} To find the total number of paths, observe that all paths will have $10$ total steps. We have to choose which $5$ of these steps will be NE (the rest will be NW). So the total number of paths is $\binom{10}{5}$. The formula for \href{https://artofproblemsolving.com/wiki/index.php?title=Combination}{combinations} is: $\binom{n}{r} = \frac{n!}{r!(n-r)!}$ and $\binom{10}{5} = \frac{10!}{5!\times5!}=252$.

 \textbf{Step 2:} Each path splits the total area of $25$ in two parts. So, for any path that gives area = $A$, you can find a unique "\textit{sister}" path that has an area = $25-A$ (in other words, the pair of paths have a combined area of 25).  Possible ways to define the "\textit{sister}" path are: 

 
                \begin{itemize}
                    \item Rotate the entire grid $180^{\circ}$
                \end{itemize}


                
                    \begin{figure}[H]
                        \centering
                        \includegraphics[width=0.5\textwidth]{Solution/AMC_8_AMC_2025_25_Solution_Figure_1.png}
                    \end{figure}
                    

 
                \begin{itemize}
                    \item Swap each step of the original paths (for example, each NW becomes NE) (this is a reflection over the diagonal)
                \end{itemize}


                \textbf{Step 3:} There are a few ways to get from this observation to the total area: 

 
                \begin{itemize}
                    \item There are $252/2 = 126$ pairs of such paths, and the total area of each pair is $25$.  So the total area given by all paths is $126 \times 25$.\item Each of the $252$ paths gives an area of $25$ if you also count the "\textit{sister}" paths. Since each "\textit{sister}" path is also one of the $252$, you have to divide by $2$ to avoid double counting.  So the total area given by all paths is $\frac{252 \times 25}{2}$.\item Note that the average area of two "\textit{sister}" paths is $\frac{25}{2}$, so you can think about every path having this area \textit{on average}.  So the total area given by all paths is $252 \times \frac{25}{2}$.
                \end{itemize}


                The final answer is $\boxed{\textbf{(B)}~3150}.$

 \textbf{Note:} This problem has a  bijection (or 1-1 correspondence) , check out \href{https://artofproblemsolving.com/store/book/intermediate-counting}{Intermediate Counting & Probability, Chapter 4}, and \href{https://artofproblemsolving.com//artofproblemsolving.com/store/book/intro-counting}{Introduction to Counting & Probability, Chapter 5}

 



\textbf{Solution 2}

If we test this problem on a smaller $2 \times 2$ diamond, we have $6$ ways to go from $A$ to $B$, and the total area is $0 + 1 + 2 + 2 + 3 + 4 = 12$, so the average area is $\frac{12}{6} = 2$, which is also the area of the diamond $2 \times 2 = 4$ divided by 2. If we assume this is true for a $5 \times 5$ diamond, then the average area is $\frac{25}{2}$. The number of paths from $A$ to $B$ is $\binom{10}{5} = 252$, and $252 \cdot \frac{25}{2} = \boxed{\textbf{(B)}~3150}$.

 This is technically flawed, since the area of a diamond is length times width.

 



\textbf{Solution 3 Easier to Motivate}

Most other solutions don't explain how they got all the cases, as well as require an insight that's somewhat hard to think of, so I'll explain another in detail. 

 If we rotate the grid $45$ degrees clockwise we can have a 5x5 grid where we can move up and right. There can only be one horizontal line segment in each column at a certain y coordinate. We'll denote the y coordinates as $0,1,2,3,4,5$. Once we choose the y coordinates for each column, we have a unique path. However, we can't move down which means that columns cannot have a higher y coordinate than the ones to the right. This is the same as distributing 5 balls in six boxes labeled $0,1,2,3,4,5$. For example, if we get 2 in 0, 2 in 1, and 1 in 4, then the order would be $0,0,1,1,4$. This is one unique path, and the total number of paths is represented by 6+5-1 choose 6-1 = 252. Also, the sum of the y-coordinates represents the area, which means we want the average sum of the y-coordinates. This is $5*(0+5)/2 = 25/2$, and $25/2 * 252 = \fbox{3150}$.

 



\textbf{Solution 4}

As found, there are $252$ total paths. However, we can take advantage of symmetry here. We consider the total area of the left side of the grid in each of the $252$ paths. Since the total area of the left side of the grid in each of the $252$ paths is equal to the total area of the right side of the grid in each of the $252$ paths, and the two total areas sum to $(252)(25) = 6300$, then the total area of the right side of the grid is $6300/2 = \fbox{3150}$.

 



\textbf{Solution 5 (maybe cheese)}

We notice that there are $\binom{10}{5} = 252$ different paths on the grid. We also notice that the average area enclosed  by a path is $\frac{25}{2}$, the area of the square divided by $2$. This means that the answer is $252 \cdot \frac{25}{2} = \fbox{3150}$.

 



\textbf{Solution 6 (q-analog)}

q-analog basic definitions are for $k\in\mathbb{N}$, then 
 \[[k]_q=1+q+\cdots+q^{k-1}\] with 
 \[[0]_q=1\] The q-factorial of $n\in\mathbb{N}$ is defined as 
 \[[n]!_q=[n]_q\cdot[n-1]_q\cdots[0]_q\] The q-binomial theorem states 
 \[\binom{n}{k}_q=\frac{[n]!_q}{[k]!_q[n-k]!_q}\] and finally for q-analog of lattice paths aka q-lattice path, we have the q-lattice path theorem which states for any $a, b\in\mathbb{N}$, 
 \[\sum_{P\in L(a,b)}q^{\text{area}(P)}=\binom{a+b}{a}_q\] To get the sum of the areas we have to differentiate the q-Lattice Path and evaluate at $q=1$. For this problem our lattice path is $L(5,5)$ and we calculate the $q$-binomial as $q^{25}+q^{24}+2q^{23}+3q^{22}+5q^{21}+7q^{20}+9q^{19}+\cdots$ and take the derivative $\left.\frac{d}{dq}\binom{10}{5}_q\right|_{q=1}$ to get $25q^{24}+24q^{23}+46q^{22}+66q^{21}+105q^{20}+140q^{19}+171q^{18}+\cdots$ and then plug in the 1 to finally get $\fbox{3150}$. ~Lopkiloinm

 shoutout to \href{https://enumeration.ca/}{https://enumeration.ca/}

 



\textbf{Solution 7 (Rigorous Representation Theoretical)}

Let $\mathcal P_{a,b}$ be the set of lattice paths from $(0,0)$ to $(a,b)$using north steps $N=(0,1)$ and east steps $E=(1,0)$.There are
 \[|\mathcal P_{a,b}|=\binom{a+b}{a}\]such paths.

 For a path $P\in\mathcal P_{a,b}$ list its steps as the word$\mathbf w=w_1\cdots w_{a+b}$ with each $w_i\in\{E,N\}$.Define the lattice area
 \[\operatorname{area}(P)=\sum_{i=1}^{a}y_i(P),\]where $y_i(P)$ is the $y$-coordinate just before the $i$-th $E$.Equivalently,
 \[\operatorname{area}(\mathbf w)= \iota(\mathbf w):= \bigl|\{(i,j):i<j,\;w_i=E,\;w_j=N\}\bigr|.\]

 The symmetric group $S_{a+b}$ acts by permuting positions of the word, and the action is transitive; in fact
 \[\mathcal P_{a,b}\cong S_{a+b}\big/\!\bigl(S_a\times S_b\bigr), \qquad M:=\mathbb C[\mathcal P_{a,b}]\cong \operatorname{Ind}^{S_{a+b}}_{S_a\times S_b}\mathbf 1.\]

 Fix an ordered pair of positions $(i,j)$ with $i<j$.Among all words, exactly half have $(w_i,w_j)=(E,N)$ and half have $(N,E)$.Hence
 \[\sum_{\mathbf w\in\mathcal P_{a,b}}\iota(\mathbf w) =\binom{a+b}{2}\,\binom{a+b-2}{a-1} =\frac{ab}{2}\,\binom{a+b}{a},\]since for each ordered pair $(i,j)$ we set $w_i=E$, $w_j=N$and distribute the remaining $a-1$ $E$’s and $b-1$ $N$’s among theother $a+b-2$ positions in $\binom{a+b-2}{a-1}$ ways.

 Therefore
 \[\boxed{\; \sum_{P\in\mathcal P_{a,b}}\operatorname{area}(P)       =\binom{a+b}{a}\,\frac{ab}{2}}\]and the average lattice area under such paths is $ab/2$.

 



\textbf{Solution 8 (bash)}

Find every single case and add them up, getting $\fbox{3150}$

 It's an annoying, but it's better than doing nothing right?(Yes I did this in-contest)(No you didnt)

 


\end{enumerate}